\section{Normal and general equation of a plane}

The normal form of a plane is one of the most useful descriptions in analytic geometry. It uses a point on the plane and a vector perpendicular to the plane.

Let
\[
P_0=(x_0,y_0,z_0)
\]
be a point on the plane, and let
\[
n=(A,B,C)\ne(0,0,0)
\]
be a normal vector. A point
\[
P=(x,y,z)
\]
lies on the plane exactly when \(P-P_0\) is perpendicular to \(n\). Therefore
\[
n\cdot(P-P_0)=0.
\]

\begin{definitionbox}
A plane through \(P_0=(x_0,y_0,z_0)\) with normal vector \(n=(A,B,C)\ne(0,0,0)\) can be written as
\[
A(x-x_0)+B(y-y_0)+C(z-z_0)=0.
\]
This is the normal equation of the plane.
\end{definitionbox}

The equation is not a memorized pattern. It says that the displacement from the fixed point to any point of the plane must be perpendicular to the normal vector.

\subsection*{From normal form to general form}

Expanding the normal equation gives
\[
Ax+By+Cz+D=0,
\]
where
\[
D=-Ax_0-By_0-Cz_0.
\]
This is called the general equation of a plane.

\begin{definitionbox}
The general equation of a plane in space is
\[
Ax+By+Cz+D=0,
\]
where \((A,B,C)\ne(0,0,0)\). The vector
\[
n=(A,B,C)
\]
is normal to the plane.
\end{definitionbox}

The coefficients of \(x\), \(y\), and \(z\) are not only numbers in an equation. Together they form a normal vector.

\subsection*{Example from a point and normal vector}

\begin{examplebox}
Find the equation of the plane passing through
\[
A=(1,0,2)
\]
with normal vector
\[
n=(2,-1,1).
\]
Use the normal form:
\[
2(x-1)-1(y-0)+1(z-2)=0.
\]
Thus
\[
2(x-1)-y+z-2=0.
\]
Expand:
\[
2x-2-y+z-2=0,
\]
so
\[
2x-y+z-4=0.
\]
\end{examplebox}

We can check the result. The point \((1,0,2)\) satisfies the equation:
\[
2\cdot1-0+2-4=0.
\]
The normal vector visible in the final equation is \((2,-1,1)\), as required.

\subsection*{Reading information from an equation}

If a plane is given by
\[
3x-2y+5z-7=0,
\]
then a normal vector is immediately visible:
\[
n=(3,-2,5).
\]
To find a point on the plane, choose convenient values for two variables and solve for the third.

\begin{examplebox}
Find one point on
\[
3x-2y+5z-7=0.
\]
Choose \(y=0\) and \(z=0\). Then
\[
3x-7=0,
\]
so
\[
x=\frac73.
\]
Thus
\[
\left(\frac73,0,0\right)
\]
lies on the plane.
\end{examplebox}

There are infinitely many possible points on a plane. We choose one that is easy to compute.

\begin{summarybox}
When reading a plane equation \(Ax+By+Cz+D=0\), ask:
\begin{itemize}
\item Are \(A\), \(B\), and \(C\) not all zero?
\item What normal vector is visible?
\item Can I find a convenient point on the plane?
\item Does a given point satisfy the equation?
\item Would the normal vector help with distance, angle, or parallelism questions?
\end{itemize}
\end{summarybox}

The general equation is powerful because it stores the normal direction directly in its coefficients.